\documentclass{article}
\usepackage{amsmath,amssymb,xurl}
\usepackage[hidelinks]{hyperref}
% Shared commands for notes in docs/paper-gaps/.

\DeclareUnicodeCharacter{2080}{\ifmmode{}_{0}\else\textsubscript{0}\fi}
\DeclareUnicodeCharacter{2081}{\ifmmode{}_{1}\else\textsubscript{1}\fi}
\DeclareUnicodeCharacter{2082}{\ifmmode{}_{2}\else\textsubscript{2}\fi}
\DeclareUnicodeCharacter{2099}{\ifmmode{}_{n}\else\textsubscript{n}\fi}

\newcommand{\Lean}{\textsc{Lean}}
\ExplSyntaxOn
\NewDocumentCommand{\leanid}{m}{
  \ifmmode
    \textsf{\detokenize{#1}}
  \else
    \group_begin:
    \urlstyle{sf}
    \tl_set:Nn \l_tmpa_tl {#1}
    \tl_replace_all:Nnn \l_tmpa_tl {\_} {_}
    \exp_args:NV \path \l_tmpa_tl
    \group_end:
  \fi
}
\ExplSyntaxOff
\providecommand{\tr}{\operatorname{tr}}
\newcommand{\tnleanIssueBase}{https://github.com/LionSR/TNLean/issues/}
\newcommand{\tnleanPrBase}{https://github.com/LionSR/TNLean/pull/}
\newcommand{\tnleanDocBase}{https://sirui-lu.com/TNLean/docs/}
\newcommand{\ghissue}[1]{\href{\tnleanIssueBase#1}{GitHub issue \##1}}
\newcommand{\ghpr}[1]{\href{\tnleanPrBase#1}{PR \##1}}
\newcommand{\doclink}[1]{\href{\tnleanDocBase#1.html}{\path{#1}}}

% Dirac notation and the identity operator, as in blueprint/src/macros/common.tex.
% \providecommand keeps note-local definitions authoritative.
\usepackage{dsfont}
\providecommand{\ket}[1]{|#1\rangle}
\providecommand{\bra}[1]{\langle #1|}
\providecommand{\braket}[2]{\langle #1 | #2 \rangle}
\providecommand{\ketbra}[2]{|#1\rangle\!\langle #2|}
\providecommand{\Id}{\mathds{1}}
\providecommand{\one}{\mathds{1}}
\providecommand{\C}{\mathbb{C}}
\providecommand{\MN}[1]{M_{#1}(\C)}
\providecommand{\MD}{M_D(\C)}

% Verdict marker: \gapnote{<kind>}{<status>}. Kind is the policy
% classification (clarification, local-correction, scope-restriction,
% unfaithful, false-source, open-gap); status is open, wip (elimination
% underway), resolved, or historical. Machine-read by the site index;
% prints a verdict line.
\newcommand{\gapnote}[2]{\par\noindent\textbf{Verdict:} #1 (#2).\par}

\title{Reduced Native Words for Empty Square Support}
\author{}
\date{}
\begin{document}
\maketitle
\gapnote{scope-restriction}{open}

Schuch, Cirac, and P\'erez-Garc\'ia, \cite[Lemma~6.14]{Schuch2010PEPS},
source \path{Papers/1001.3807/paper_v3.tex}, lines 2199--2214,
deform a fixed-endpoint flux string by virtual invariance.
The result here concerns the generated path relation in the particular case
where no primal square support is allowed. It does not assert the source's
unrestricted geometric deformation claim.

Let $X=\mathbb Z_w\times\mathbb Z_h$ with $w,h>0$. A labelled native step
has a positive east or north edge, based at an actual vertex of $X$, and
a sign. West and south steps use inverse letters of those same edges.
For a finite path $p$, let $L(p)$ be its signed edge list in reverse traversal
order. Cancelling adjacent inverse letters gives the reduced list
$\operatorname{red}(p)$. This retains the actual base vertices and directions,
including the separately labelled edges of periods one and two.

Write $p\sim_{\varnothing}q$ for the generated relation allowing backtracks,
concatenation and reversal, with no square moves. For parallel paths,
\begin{align}
    p\sim_{\varnothing}q
    \quad\Longleftrightarrow\quad
    \operatorname{red}(p)=\operatorname{red}(q).
    \label{eq:scp10_empty_support_normal_form}
\end{align}
The forward direction evaluates the constant free-group transport
$F_e(z)=W(e)z$ at the identity, where $W(p)$ is the free-group word represented
by $L(p)$. Thus $F_p(1)=W(p)$ and a generated homotopy forces equality of words.
Conversely, reduced-word equality implies equality in the free group,
then equality in the finite native presentation, whose completeness supplies
the generated homotopy. In fact this converse works for every allowed
support set; necessity of reduced-word equality here uses empty support.
\footnote{\leanid{TNLean.PEPS.torusDualHomotopy_of_reducedWord_eq};
\leanid{TNLean.PEPS.torusDualHomotopy_empty_iff_reducedWord_eq}.}

The actual Boolean procedure reduces the two finite signed lists and compares
the results:
\begin{align}
    D(p,q)=[\,\operatorname{red}(p)=\operatorname{red}(q)\,],\qquad
    D(p,q)=\mathrm{true}\quad\Longleftrightarrow\quad p\sim_{\varnothing}q.
    \label{eq:scp10_empty_support_decision}
\end{align}
Its computation uses only finite-list reduction and equality of native letters,
with no decision oracle for a presented-group word problem.
\footnote{\leanid{TNLean.PEPS.torusDualStepLetter};
\leanid{TNLean.PEPS.torusDualPathLetters};
\leanid{TNLean.PEPS.torusDualReducedWord};
\leanid{TNLean.PEPS.torusDualEmptyHomotopyDecision};
\leanid{TNLean.PEPS.torusDualEmptyHomotopyDecision_eq_true}.}

The general proper nonempty support still requires a geometric normal form
for the allowed face relations. The finite-presentation criterion alone
does not supply that algorithm. This result supplies no physical density
classification and does not identify distinct winding or endpoint sectors.
The nonabelian restricted-homotopy and physical annular obstructions remain
separate. The broader work stays under \ghissue{8703}; \ghissue{8947} is unchanged.

\bibliographystyle{alpha}
\bibliography{references}
\end{document}
